\section{Protocol Design}
\label{sec:protocol}

We implement BPE on Ethereum (Base L2) using Superfluid's General Distribution
Agreement (GDA)~\cite{superfluid2024}. The architecture comprises four contracts.

\subsection{Architecture Overview}

\begin{figure}[ht]
\centering
\begin{tikzpicture}[
  box/.style={draw, rounded corners, minimum width=2.8cm, minimum height=0.9cm, align=center, font=\small},
  arr/.style={-{Stealth[length=3mm]}, thick},
  lbl/.style={font=\scriptsize, midway, above},
  >=Stealth,
]
  % Nodes
  \node[box] (sm)  at (0,0)   {StakeManager\\[-2pt]\scriptsize stake, $\sqrt{\cdot}$ cap};
  \node[box] (cr)  at (4.5,0) {CapacityRegistry\\[-2pt]\scriptsize EWMA, commit-reveal};
  \node[box] (bp)  at (4.5,-2){BackpressurePool\\[-2pt]\scriptsize rebalance};
  \node[box] (gda) at (9,-2)  {Superfluid GDA\\[-2pt]\scriptsize streaming dist.};
  \node[box] (eb)  at (1.5,-4){EscrowBuffer\\[-2pt]\scriptsize overflow hold};
  \node[box] (pl)  at (6,-4)  {Pipeline\\[-2pt]\scriptsize multi-stage};
  \node[box] (pc)  at (9,0)   {PricingCurve\\[-2pt]\scriptsize dynamic fees};
  \node[box] (ct)  at (0,-2)  {CompletionTracker\\[-2pt]\scriptsize verification};
  \node[box] (oa)  at (0,-4)  {OffchainAggregator\\[-2pt]\scriptsize EIP-712 batch};

  % Edges
  \draw[arr] (sm)  -- node[lbl]{cap} (cr);
  \draw[arr] (cr)  -- node[lbl,right,pos=0.4]{capacity} (bp);
  \draw[arr] (bp)  -- node[lbl]{units} (gda);
  \draw[arr] (bp)  -- (eb);
  \draw[arr] (bp)  -- (pl);
  \draw[arr] (cr)  -- (pc);
  \draw[arr] (cr)  -- (ct);
  \draw[arr] (oa)  -- (cr);
  \draw[arr, dashed] (ct) -- node[lbl,left,pos=0.5]{\scriptsize slash} (sm);
\end{tikzpicture}
\caption{BPE protocol architecture. Solid arrows denote read dependencies;
the dashed arrow represents the auto-slash feedback path from
CompletionTracker to StakeManager.}
\label{fig:architecture}
\end{figure}

\subsection{StakeManager}

Sinks deposit ERC-20 tokens to stake. The capacity cap is computed as
\begin{equation}\label{eq:capacity-cap}
\text{cap}(k) = \sqrt{\text{stake}(k) / u}
\end{equation}
where $u$ is the stake unit
parameter (set to $10^{18}$). This concave function implements Sybil resistance:
splitting stake $S$ into $n$ sinks yields total capacity
$n \cdot \sqrt{S/(nu)} = \sqrt{n} \cdot \sqrt{S/u} < n \cdot \sqrt{S/u}$
for $n > 1$---strictly less than the capacity from a single identity.
A minimum per-sink stake $S_{\min}$ further penalizes splitting.

\subsection{CapacityRegistry}

The registry manages task type creation, sink registration, and capacity
signaling. Capacity updates use a two-phase commit-reveal protocol for
MEV resistance:

\begin{enumerate}
  \item \textbf{Commit}: Sink submits $h = \texttt{keccak256}(C, \text{nonce})$.
  \item \textbf{Reveal}: Within 20 blocks, sink reveals $(C, \text{nonce})$.
        The contract verifies the hash, applies the stake cap, and computes
        the EWMA update: $\Csmooth_{\text{new}} = 0.3 \cdot C_{\text{capped}}
        + 0.7 \cdot \Csmooth_{\text{old}}$.
\end{enumerate}

The registry maintains an aggregate \texttt{totalCapacity} per task type,
updated atomically on each reveal.

\subsection{BackpressurePool}

For each task type, the pool creates a Superfluid GDA pool and acts as its
admin. On \texttt{rebalance()}, it reads all sink capacities from the registry
and sets each sink's pool units proportional to their smoothed capacity:
$\text{units}(k) = \lfloor C_{\text{smooth}}(k) \cdot 10^9 / C_{\text{total}} \rfloor$.

Sources create Constant Flow Agreement (CFA) streams \emph{to the pool address}.
The GDA automatically distributes incoming flows to members proportional to their
units---achieving capacity-proportional payment routing without centralized scheduling.

Rebalance is \emph{permissionless}: any address may call it. A 5\% threshold
(measured in basis points of total capacity change) prevents unnecessary gas
expenditure.

\subsection{EscrowBuffer}

When all sinks in a task type are at capacity, excess payments are deposited
into the buffer. The \texttt{drain()} function distributes buffered funds
proportionally to current sink capacities. A configurable \texttt{bufferMax}
caps the buffer size; when full, sources must pause.

\subsection{Pipeline Composition}

For multi-stage agent pipelines, the \texttt{Pipeline} contract chains
BackpressurePools. The key mechanism is \emph{upstream propagation}:
each stage's effective capacity is constrained by its downstream stage:
\begin{equation}
C_{\text{eff}}(\text{stage}_i) = \min\bigl(C_{\text{local}}(\text{stage}_i),\;
C_{\text{eff}}(\text{stage}_{i+1})\bigr)
\end{equation}
Rebalance proceeds back-to-front, ensuring downstream congestion propagates
upstream before any stage is rebalanced.
